\subsection{Octad}

The Octad is a short chapter.

Iamblichus continues his march\footnote{\url{http://www.scribd.com/doc/48888397/Iamblichus-the-Theology-of-Arithmetic}} through numbers noting that all men, without exception, count 1-10. The empiricist says this is because they have 10 fingers, and the Hermeticist says ``why do you suppose that is?''. In any case, it is not impossible to conceive attempts to circumvent divine order, such as the reordering of week days\footnote{\url{http://en.wikipedia.org/wiki/French_Republican_Calendar}} during the French Revolution, so their argument holds little water (which they know -- most of the Anglos don't use metric).

Eight is perfect, in that it is 1+7, neither number which is engendered, as Seven is a ``second'' Monad, in this meditation. The first two odd numbers (3 and 5) also produce it, making it relate to the process of the Cube (Claims Iamblichus, although he doesn't specifiy how this is so).

The pattern four is in the ecliptics of the Zodiac, and the ``circles'' (tropics and polar), as well as finding expression in Harmonic divisions and even the teeth and skulls\footnote{\url{http://en.wikipedia.org/wiki/Human_skull}} (presumably four big plates) of human beings. For instance, birds usually have four claws, and if an insect has more than eight legs, it is usually not exact.

\begin{quotex}
Philolaus teaches that after mathematical magnitude has become three dimensional thanks to the tetrad, there is the quality and `color' of visible Nature in the pentad, and the ensoulment in the hexad, and intelligence and health and what he calls `Light' in the hebdomad, and then next, with the ogdoad, things come by love and friendship and wisdom and creative thought.
\end{quotex}


So we are looking more now at what is normally understood as ``manifestation'' or ``Creation''. It is safety and foundation, because its root is in 2, and 2 is the leader of Manifestation in the sense of ``daring'' or ``courage'' (warrior-qualities).

There is a long section on harmonic relationships which is very involved with terms like sesquitertian, sesquioctave, etc., but I will recommend it to musicians, only noting for the Layman that what he needs to know is that the sesqui interval works as follows. 12 is the sesquialter of 8 (it succeeds 8 by an interval of half, so the ``change'' is related to the ``substance'' by law), and the sesquitertian of 9 (it succeeds 9 by an interval of one-third).

There is, of course, the Indian method\footnote{\url{http://www.sanatansociety.org/vedic_astrology_and_numerology/indian_numerology_8.htm}} of examining the birthday, name, and age of the person in order to use Pythagorean reduction to basic numerals to arrive at spiritual meaning. But these are more personal methods, and less metaphysical, than those which Iamblichus deploys.

8 resembles the infinity symbol, turned at a right angle. Does this mean that it is Infinity, set on its ear? In other words, Infinity set towards a downfall? Iamblichus does seem to indicate that ``manifestation'' in terms of Bios begins to break forth in 8, interweaving design and matter in actual genesis. We might think of the TV programs we were all indoctrinated on, which always start with some image of algae or lava or star dust, coalescing and shimmering and out-breaking, in order to appreciate the power of 8. The affinities are revealed: sex, division, mutual interdependence, etc.

Just for fun I looked up some modern\footnote{\url{http://www.youtube.com/watch?v=Adlw7sRw7cg}} numerology; Ellis Taylor seems to think that our Arabic numbers (in contradistinction to Latin ones, which are more pure) have been used to program us for failure. Although this sounds like the typical modern conspiracy theory, it may have something of interest, considering that even the Solfeggio\footnote{\url{http://www.miraclesandinspiration.com/solfeggiofrequencies.html}} scale has been reset in the modern Era. All changes of this sort ought to be more self-consciously investigated.

\begin{quotex}
Our modern day musical scale is slightly out of sync from the original Solfeggio frequencies and is, consequently, more dissonant as it is based upon what is termed the ``Twelve-Tone Equal Temperament.'' In ancient times, the musical scale was called ``Just Intonation.'' And also, our modern music falls within the A 440 hz frequency, which was changed from A 417 hz, around 1914.

In addition, a 7th note was added in the form of a ``SI,'' or a ``TI,'' as in the ``DO, RE, MI, FA, SO, LA, TI'' vocal scale, while the original Solfeggio scale was composed of only six notes: ``UT, RE, MI, FA, SO, LA.''
\end{quotex}


I am not a literal conspiracist, but certainly there was a reason the musical scale was altered: Traditionalists who are mathematicians or musicians ought to look into these things. We recall that Gornahoor has taught that there is an invisible war on a higher pattern, and that changes are never random, even if no discernible pattern can be found.


\flrightit{Posted on 2013-03-29 by Logres}

\begin{center}* * *\end{center}

\begin{footnotesize}
\begin{sffamily}
\texttt{George Goerlich on 2013-03-31 at 23:49 said:}

Sorry for the ignorance, but are there examples of recorded ancient/traditionalist music?

\hfill

\texttt{Logres on 2013-04-01 at 11:13 said:}

Plato discusses music (eg., his approval of certain ``modes'' such as Dorian or Phrygian); I've noted the Solfeggio shift that ought to be investigated. The modern mania for ``kid's music'' (rock and roll) and other even more degenerate forms (it's hard not to conclude that jazz is really such a form), even when better is available, is notable. I am not a musician, but a lot of Celtic music is played (I think) in an off mode that is either Dorian or Phrygian -- there really aren't any accidents. Musicians either place us closer to God or closer to the devil; traditionalists can use their souls to verify which is which, even if they don't have a lot of theory to go on. Augustine, I believe, has a very useful discourse on music. Listen to some of the Templar chants, like Salve Regina or The Stone that the Builders Rejected.

\hfill

\texttt{Arthur on 2013-04-01 at 22:52 said:}

Conspiracy or not, here is something about the imposing of the 440 herz pattern:

\url{http://www.medicalveritas.org/MedicalVeritas/Musical_Cult_Control.html}


\end{sffamily}
\end{footnotesize}

